\scanpage{308}
\mspage{454}

\uline{V}\qquad
\uline{Study}\inserted{\uline{(of the equation)}}\uline{ in}
$GP(1)$, \uline{and the condition} $\mathrm{Tr}\,\sigma=0$.

Up to now we have studied equation (32)
\[
\tag{*}
[\alpha,\rho]=[\beta,\unsure{\sigma}]\,
\unsure{\varepsilon_1(\mu-1)}
\]
in $G\deleted{\unsure{\ell_2}}$ ($\simeq\mathrm{Gl}(2)$), or what amounts to the
same, \deleted{in} as a relation in $\mathrm{Sl}(2)$.
Here $\alpha,\beta$ are \deleted{\illegible{1}}\inserted{variable} sections of
$G\deleted{\unsure{\ell_2}}\simeq\mathrm{Gl}(2)$, but $[\alpha,\rho]$ and
$[\beta,\sigma]$ depend only on the \unsure{corresponding} sections
$\widetilde\alpha,\widetilde\beta$ of
\[
PG\ (\simeq GP(1))\simeq
G/\underbrace{\unsure{G_m}}_{\text{Center of }G}
\simeq\mathrm{Aut}_{\unsure{1|l}}(G)
\]
(likewise they depend on $\rho,\sigma$ only
\unsure{via} $\widetilde\rho,\widetilde\sigma$).
We can therefore \unsure{consider} that it is an equation
\inserted{in \deleted{$\mathrm{Sl}(2)$} $SG\simeq\mathrm{Sl}(2)$,}
relating $\widetilde\alpha,\widetilde\beta,\mu$.

We can propose to study the corresponding equation in
$\unsure{PG}\simeq GP(1)$,
\[
\tag{1}
[\widetilde\alpha,\widetilde\rho]
 =[\widetilde\beta,\widetilde\sigma]\cdot\widetilde\varepsilon_1(\mu-1)
\]
which can also be written in $\mathrm{Gl}(2)$, in the \unsure{equivalent} form
\[
\tag{2}
[\alpha,\rho]=\varepsilon[\beta,\sigma]\,
\overset{\deleted{\unsure{N_0}}}{\varepsilon_1}(\mu-1),
\]
where $\varepsilon$ is a section of $\unsure{G_m}$,
\deleted{and in}\inserted{of} \deleted{fact} course uniquely
\unsure{determined|determined} by $\alpha,\beta,\mu$,
(and even by $\widetilde\alpha,\widetilde\beta,\mu$)).
We will necessarily have

\scanpage{309}

\[
\tag{3}
\varepsilon^2=1
\]
(take the determinant of the two \unsure{sides}). We regard (2) as an equation
relating the variables $\alpha,\beta,\mu,\varepsilon$
($\alpha,\beta$ in $\mathrm{Gl}(2)$, $\mu,\varepsilon$ in
$\unsure{\mathbb E^1|\mathbb E'}$ -- and in fact necessarily in
$\unsure{\mathbb G_m|G_m}$, and $\varepsilon$ \unsure{already} in
$\mu_2={}_2\unsure{\mathbb G_m|G_m}$), or also
$\widetilde\alpha,\widetilde\beta,\mu,\varepsilon$.
Let us put, \deleted{as in}\inserted{\unsure{by analogy with}} (II),
\[
\tag{4}
\zeta=[\alpha,\rho],\qquad
\eta=\deleted{\illegible{1}}\varepsilon
[\unsure{\uuline{\beta}},\sigma]\deleted{\unsure{\varepsilon_1 t|\text{and}}},
\]
so that (2) takes the form
\[
\tag{5}
\zeta=\eta\varepsilon_1(\mu-1)
\]
as in (II). To be able to apply (II) to this situation, we will have to
assume that $\eta=\varepsilon[\beta,\sigma]$ satisfies the
same preliminary conditions as
\[
\eta_0=[\beta,\sigma]
\]
itself, namely $\det\eta=1$
($\Longleftrightarrow\ \varepsilon^2=1$) and
\deleted{\unsure{that}}
\[
\underbrace{\mathrm{Tr}\,\eta}_{\inserted{\varepsilon\mathrm{Tr}\,\eta_0
=\varepsilon\mathrm{Tr}\,\sigma}}
=\mathrm{Tr}\,\sigma
\qquad\deleted{\unsure{for}\ \varepsilon\mathrm{Tr}\,\eta_0=}
\]
i.e.
\[
\varepsilon\mathrm{Tr}\,\sigma=\mathrm{Tr}\,\sigma.
\]

\scanpage{310}
\mspage{455}

\uline{Lemma.} Let $\sigma$ be a section of $G$ (a form of
$\mathrm{Gl}(2)$, \unsure{associated|associated} with an Azumaya algebra
$\unsure{\mathcal M|M}$) -- i.e. a section of
$\unsure{\mathcal M|M}$ such that $\det\sigma$ invertible.
\deleted{Equivalent conditions}\inserted{Let $A_\sigma$ be the scheme}
of sections $\eta$ \deleted{\illegible{1}} of $G$ satisfying the conditions
\[
\det\eta=1\quad\text{(i.e. $\eta$ in $SG$)},\qquad
\mathrm{Tr}\,\eta\circ\sigma=\mathrm{Tr}\,\sigma.
\]
Equivalent conditions a) b) c) :

a) $A_\sigma$ is stable under translations by
$\mu_2\subset\unsure{\mathbb G_m|G_m}$ :
\[
\unsure{\eta}\longmapsto\varepsilon\eta\qquad(\varepsilon^2=1)
\]

b) $\mathrm{Tr}\,\sigma=0$

\deleted{c) $\sigma^2\in\Gamma\unsure{\mathbb G_m|G_m}$ i.e.
$\widetilde\sigma^2=1$.}

\unsure{d)|c)} (Under the \unsure{conditions} of the preceding number
\inserted{\uline{$\overline{\mathrm{IV}}$} (page 452)},
with the normalizations (34) (38) for $\rho,\sigma$) We have
\[
\tag{6}
\sigma=\begin{pmatrix}0&-1\\v'&0\end{pmatrix}
=\begin{pmatrix}0&-1\\-v&0\end{pmatrix}
\]
i.e. (with the notation of loc. cit.) we have $w=1$,
\[
\tag{7}
\text{i.e.}\qquad
\rho=\begin{pmatrix}0&1\\v&1\end{pmatrix}.
\]

\inserted{d) We have $\sigma^2\in\Gamma\unsure{\mathbb G_m|G_m}$,
i.e. $\widetilde\sigma^2=1$}

\inserted{Finally, these conditions imply the following d), and
are equivalent to \unsure{it|them} \unsure{if $\sigma$}
\deleted{\unsure{when $\sigma$}} is not a homothety on
any \unsure{fiber} (\unsure{thus} i.e. $\sigma$ is regular in the sense of
Steinberg).}

\deleted{This means that}\inserted{\unsure{The equivalence}
$a)\Longleftrightarrow b)$ means \unsure{also|thus}}
that $\unsure{\mathbb E^1|\mathbb E'}{}^{\mu_2}=(0)$, which is well
\unsure{known} and immediate to verify.
The equivalence of b) and d) is immediate, remains to prove
$b)\Longleftrightarrow c)$.
\deleted{\unsure{Let us assume|Assuming} this condition \unsure{satisfied}}
\inserted{Now we have the equation \unsure{of|of}
\unsure{Hamilt-Cayley|Hamilton-Cayley}}
\[
\sigma^2-(\mathrm{Tr}\,\sigma)\sigma+\det\sigma=0
\]

\scanpage{311}

which shows that $\mathrm{Tr}\,\sigma=0\Longrightarrow\sigma=\det\sigma\cdot1$,
and conversely if $\sigma^2$ is a scalar, we will have that
$(\mathrm{Tr}\,\sigma)\sigma=\sigma^2-\det\sigma\cdot1$ is one,
which is possible (if $\sigma$ is \deleted{\unsure{not}} a scalar on no
fiber) only if $\mathrm{Tr}\,\sigma=0$.

Let us assume these conditions satisfied for what follows,
\deleted{we \unsure{find} then that the scheme $\mathcal G_{\rho,\sigma}$
of \illegible{1} $(\alpha,\beta,\mu,\varepsilon)$ satisfying}
\inserted{Let}
\[
\tag{8}
\mathcal G^*_{\rho,\sigma}\subset
G\times G\times\unsure{\mathbb E^1|\mathbb E'}
\times\unsure{\mathbb E^1|\mathbb E'}
\]
be the scheme of solutions in $(\alpha,\beta,\mu,\varepsilon)$ of equation (2),
which is also the inverse image in
$G\times G\times\unsure{\mathbb E^1|\mathbb E'}
\times\unsure{\mathbb E^1|\mathbb E'}$ \unsure{of a} subscheme
\deleted{$\widetilde{\mathcal G}_{\rho,\unsure{\sigma}}$
\illegible{1} $PG\times$}
\[
\tag{9}
\widetilde{\mathcal G}^*_{\rho,\sigma}\subset
PG\times PG\times\unsure{\mathbb E^1|\mathbb E'}
\deleted{\times\unsure{\mathbb E^1|\mathbb E'}}
\]
consisting of the solutions \inserted{$(\alpha,\beta,\mu)$} of (1).
The projection onto the \unsure{second|last} factor
$\mathbb E^1$ gives a morphism
\[
\tag{10}
\mathcal G^*_{\rho,\sigma}\xrightarrow{\underline{\varepsilon}}\mu_2,
\]

\scanpage{312}
\mspage{456}
which moreover \unsure{factors} through $\widetilde{\mathcal G}^{*}_{\rho,\sigma}$ :
\[
\tag{11}
\widetilde{\varepsilon}:\widetilde{\mathcal G}^{*}_{\rho,\sigma}\longrightarrow\mu_2.
\]
Let us note that we have
\[
\deleted{\unsure{(11)}}\qquad
\mathcal G_{\rho,\sigma}\subset\mathcal G^{*}_{\rho,\sigma}
\]
more precisely
\[
\tag{12}
\mathcal G_{\rho,\sigma}=\underline{\varepsilon}^{-1}(1)
\]
which leads us to introduce also
\[
\tag{13}
\underset{\substack{\Vert\ \text{definition}\\
\underline{\widetilde{\varepsilon}}^{-1}(1)}}{\widetilde{\mathcal G}_{\rho,\sigma}}
=\widetilde{\mathcal G}^{*}_{\rho,\sigma},
\]
so that $\mathcal G_{\rho,\sigma}$ is the inverse image in
$G\times G\times\unsure{\mathbb E^1|\mathbb E'}$ of $\widetilde{\mathcal G}_{\rho,\sigma}$.

We have a morphism
\[
\tag{14}
\left\{\begin{aligned}
(\alpha,\beta,\mu,\varepsilon)&\longrightarrow
\left(\underbrace{(\zeta=\alpha\rho(\alpha)^{-1},\,
\eta=\varepsilon\beta\sigma(\beta)^{-1},\,\mu)}_{\in\Gamma\mathcal X_{\rho,\sigma}},\,
\underbrace{\varepsilon}_{\in\Gamma\mu_2}\right)\\
\text{from}\quad\mathcal G^{*}_{\rho,\sigma}&\longrightarrow\mathcal X_{\rho,\sigma}\times\mu_2
\end{aligned}\right.
\]
and the inverse image of a section $((\zeta,\eta,\mu),\varepsilon)$ of the right-hand side is the
scheme of solutions in $\alpha,\beta$ of the equations
\[
\tag{15}
{}^{(a)}\underset{\inserted{\unsure{\alpha}}}{\alpha\rho(\alpha)^{-1}}=\zeta,
\qquad {}^{(b)}\beta\sigma(\beta)^{-1}=\varepsilon\eta
\]

\scanpage{313}
\unsure{Now|We} we have, via the isomorphism
\[
B'_0\xrightarrow{\sim}\mathcal X_{\rho,\sigma}
\]
a privileged solution $\alpha=\underline u$ of the first equation (15), which is
\unsure{characterized} as the unique solution which is \deleted{\unsure{a}} section of
$B'_0$ (\unsure{uniqueness} resulting from
\[
T_\alpha\cap B'_0=\{1\}
\]
).
But we \unsure{will have}, for this $u$
\[
\underbrace{u\sigma(u)^{-1}}_{[u,\sigma]}=\eta
\qquad(\text{and not }\varepsilon\eta\ !)
\]
to find in terms of $u$ a solution of the equation
\inserted{(15 b)} $\beta\sigma(\beta)^{-1}=\varepsilon\eta$, we are led to try in the form
\[
\tag{16}
\beta=u\tau\quad\text{whence}\quad
\beta\deleted{\unsure{\rho}}\inserted{\sigma}(\beta)^{-1}
=\deleted{\unsure{\beta}}\inserted{u}\tau\sigma(\tau)^{-1}
\sigma(\deleted{\unsure{\beta}}\inserted{u})^{-1},
\]
trying to determine $\tau$ so that we have
\[
\tag{17}
\underbrace{\tau\sigma(\tau)^{-1}}_{[\tau,\sigma]}=\varepsilon,
\]
which indeed implies
\[
\beta\sigma(\beta)^{-1}
=\deleted{\unsure{\beta}}\inserted{u}\varepsilon\sigma(\deleted{\unsure{\beta}}\inserted{u}^{-1})
=\varepsilon\deleted{\unsure{\beta}}\inserted{u}\sigma(\deleted{\unsure{\beta}}\inserted{u})^{-1}
=\varepsilon\eta
\]
i.e. (15 b). \deleted{\unsure{Lemma}}

\scanpage{314}
\mspage{457}
The relation (17) can also be written
\[
\deleted{\sigma(\tau)=\varepsilon\tau}
\qquad(\text{take the inverse})
\]
\[
\tag{18}
\tau\sigma\tau^{-1}\sigma^{-1}=\varepsilon
\quad\text{or}\quad
\sigma\tau\sigma^{-1}\tau^{-1}=\varepsilon
\quad\text{or}
\]
\[
\tag{19}
\sigma=\varepsilon\tau(\sigma)
\quad\text{or}\quad
\tau(\sigma)=\varepsilon\sigma
\]
which implies that $\tau$ \uline{normalizes} $T_\sigma$
\[
\tag{20}
\tau(T_\sigma)=T_\sigma\quad\text{i.e.}\quad\tau\in\Gamma N_\sigma
\]
where
\[
\tag{21}
N_\sigma\overset{\text{definition}}{=}\operatorname{Norm}_G(T_\sigma).
\]

\uline{Corollary to the lemma} (page 455) the \inserted{\unsure{equivalent}} conditions of the lemma
\deleted{\unsure{are equivalent}} \inserted{are implied}
\deleted{to the} \inserted{by the} following ;
\inserted{and \unsure{imply it} if $\sigma$ is \unsure{regular|regular} :}

(e) $\forall$ \inserted{fppf local} section
$\varepsilon$ of $\mu_2$, exists \inserted{\unsure{local\textsuperscript{ly}|locally}}
section \deleted{\illegible{1}} $\tau$ of $G$ such that we have (17) (or again (19))

(f) (If $2$ invertible on $S$) idem, with $\varepsilon=-1$.

If we have (e), then we see that for every \deleted{\illegible{1}} $\varepsilon$,
we must have $\operatorname{Tr}\varepsilon\sigma=\operatorname{Tr}\sigma$
now $\operatorname{Tr}\varepsilon\sigma=\varepsilon\operatorname{Tr}\sigma$, hence
\inserted{(varying $\varepsilon$)}
$\operatorname{Tr}\sigma\in\Gamma(\unsure{\mathbb E^1|\mathbb E'}^{\mu_2}=0)$
hence $\operatorname{Tr}\sigma=0$. If we assume it only for $\varepsilon=-1$,
we will have $\operatorname{Tr}\sigma=-\operatorname{Tr}\sigma$ i.e.
$2\operatorname{Tr}\sigma=0$, i.e. $\operatorname{Tr}\sigma=0$ if $2$ invertible.
Conversely, let us assume $\operatorname{Tr}\sigma=0$, then
$\operatorname{Tr}\varepsilon\sigma=0$, moreover
$\det(\varepsilon\deleted{\unsure{\sigma}})=\varepsilon^2=1$, hence $\varepsilon$ satisfies
the ``numerical conditions'' to be of the form $[\tau,\sigma]$ for
suitable $\tau$. For